%! Sine, Cosine, and Tangent — Unit 3, Lesson 3.2 — Warm-Up (Answer Key)
\documentclass[10pt]{article}
\usepackage{precalculus-article}
\usepackage{precalculus-key}   % pulls in -boxes; do NOT also load -boxes
\newcommand{\TallMath}[1]{$\displaystyle #1\rule[-1.4em]{0pt}{3.2em}$}

\begin{document}

\pageheader{Unit 3, Lesson 3.2}{Warm-Up --- Answer Key}
\namedateperiod

\vspace{0.10in}

\begin{notesbox}{Spiral Review: Coordinates, Distance, and Angle Rotation --- Key}

\textbf{1.} A point is located at coordinates $(3,\, 4)$. What is its distance from the origin?
Show your work using the Pythagorean theorem.

\medskip
\ans{$d = \sqrt{3^2 + 4^2} = \sqrt{9 + 16} = \sqrt{25} = 5$}

\vspace{0.18in}

\textbf{2.} A point $P = (0.6,\; 0.8)$ lies on a circle of radius 1.
\begin{itemize}[leftmargin=1.6em, itemsep=2pt]
  \item[(a)] Verify that $P$ is on the unit circle by computing $x^2 + y^2$. \hfill
        \ans{$(0.6)^2 + (0.8)^2 = 0.36 + 0.64 = 1$ \checkmark}
  \item[(b)] Compute the ratio $y/x$ for this point. \hfill
        \ans{$0.8/0.6 \approx 1.33$}
\end{itemize}

\vspace{0.18in}

\textbf{3.} An angle starts at the positive $x$-axis and rotates $90°$ counterclockwise.
Circle the direction the terminal ray points:

\medskip
\ans{positive $y$-axis}

\vspace{0.18in}

\textbf{4.} A wheel of radius 1 completes one full revolution.
\begin{itemize}[leftmargin=1.6em, itemsep=2pt]
  \item[(a)] How many degrees is one full revolution? \hfill \ans{$360°$}
  \item[(b)] The arc length traced equals the angle measured in radians when $r = 1$.
        What arc length does one full revolution trace? \hfill \ans{$2\pi$}
\end{itemize}

\end{notesbox}

\end{document}
